% =========================================================================
%  WCN Lab Assignment -- main LaTeX source
%  The USN is NOT typed here. run_all.m reads it from config.m and writes
%  generated/usn.tex; it also writes generated/expNN_code.m (the merged code
%  listing) and generated/expNN_outputs.tex (the output images).
%  Build:  pdflatex -jobname=WCN_Lab_Assignment main.tex   (run twice)
% =========================================================================
\documentclass[11pt,a4paper]{article}

\usepackage[a4paper,top=25mm,bottom=25mm,left=24mm,right=24mm,headheight=14pt]{geometry}
\usepackage[T1]{fontenc}
\usepackage{lmodern}
\usepackage[scaled=0.86]{beramono}
\usepackage{microtype}
\usepackage{graphicx}
\usepackage{float}
\usepackage{caption}
\usepackage{listings}
\usepackage{fancyhdr}
\usepackage{titlesec}
\usepackage{array}
\usepackage{chngcntr}
\usepackage[hidelinks]{hyperref}

\newcommand{\USN}{YOUR\_USN\_HERE}
\newcommand{\RunEnv}{GNU Octave 11.3.0}


% ---------------------------------------------------------------- headings
\titleformat{\section}[display]
  {\normalfont\sffamily\centering}
  {\LARGE\bfseries Experiment~\thesection}
  {6pt}
  {\Large\bfseries}
  [{\vspace{4pt}\titlerule[0.8pt]}]
\titlespacing*{\section}{0pt}{0pt}{14pt}
\titleformat{\subsection}{\normalfont\sffamily\large\bfseries}{}{0pt}{}
\titlespacing*{\subsection}{0pt}{14pt}{6pt}

% ---------------------------------------------------------------- header/footer
\pagestyle{fancy}
\fancyhf{}
\fancyhead[L]{\small\sffamily Wireless and Cellular Networks Laboratory (23ET7PCWCN)}
\fancyhead[R]{\small\sffamily USN: \USN}
\fancyfoot[C]{\small\sffamily\thepage}
\renewcommand{\headrulewidth}{0.4pt}

% ---------------------------------------------------------------- code (black only)
\lstset{
  language=Matlab,
  basicstyle=\ttfamily\small,
  keywordstyle={}, commentstyle={}, stringstyle={}, identifierstyle={},
  numbers=left, numberstyle=\tiny, numbersep=8pt,
  frame=single, framesep=4pt,
  xleftmargin=18pt, framexleftmargin=4pt,
  breaklines=true, breakatwhitespace=false,
  postbreak=\mbox{$\hookrightarrow$\space},
  columns=fullflexible, keepspaces=true,
  showstringspaces=false, tabsize=4, upquote=true,
  literate={*}{{\char42}}1 {-}{{\char45}}1
}

% ---------------------------------------------------------------- figures: 1.1, 1.2, ...
\counterwithin{figure}{section}
\captionsetup{font={small,sf},labelfont=bf,skip=6pt}
\newcommand{\consoleoutput}[1]{%
  \begin{figure}[H]\centering
  \fbox{\includegraphics[width=0.95\linewidth,height=0.78\textheight,keepaspectratio]{#1}}
  \caption*{Command Window Output}\end{figure}}
\newcommand{\figureoutput}[1]{%
  \begin{figure}[H]\centering
  \includegraphics[width=0.92\linewidth,height=0.62\textheight,keepaspectratio]{#1}
  \caption{}\end{figure}}

% ---------------------------------------------------------------- experiment block
% \experiment{NN}{title}{aim}{check value from the manual}
\newcommand{\experiment}[4]{%
  \clearpage
  \section{#2}
  \subsection{Aim}
  #3
  \par\smallskip
  {\small\textbf{Check value (lab manual):} \emph{#4}}
  \subsection{Code}
  \lstinputlisting{generated/exp#1_code.m}
  \subsection{Output}
  \input{generated/exp#1_outputs.tex}}

\begin{document}

% ======================================================================== title
\begin{titlepage}
\centering\sffamily
\vspace*{18mm}
{\large BMS COLLEGE OF ENGINEERING, BENGALURU\par}
\vspace{2mm}
{\normalsize Department of Electronics and Telecommunication Engineering\par}
\vspace{28mm}
\rule{\linewidth}{1.2pt}\par
\vspace{6mm}
{\Huge\bfseries Wireless and Cellular\\[4pt] Networks Laboratory\par}
\vspace{5mm}
{\Large Lab Assignment --- MATLAB Experiments 1 to 15\par}
\vspace{6mm}
\rule{\linewidth}{1.2pt}\par
\vspace{10mm}
{\large Course Code: 23ET7PCWCN \quad$\cdot$\quad Semester VII \quad$\cdot$\quad Academic Year 2026--27\par}
\vfill
\renewcommand{\arraystretch}{1.6}
\begin{tabular}{>{\bfseries}r l}
USN & \USN \\
Execution environment & \RunEnv \\
\end{tabular}
\vspace{22mm}
\end{titlepage}

% ======================================================================== contents
\pagenumbering{roman}
\tableofcontents
\thispagestyle{fancy}
\clearpage
\pagenumbering{arabic}

\experiment{01}{Frequency Reuse, Co-channel Interference and Trunked Capacity}
{To construct a hexagonal cellular layout, evaluate the co-channel reuse ratio and the signal to
interference ratio for admissible cluster sizes, and determine the traffic capacity per cell using
the Erlang~B model.}
{For $N = 7$ the co-channel reuse ratio is $\sqrt{21} = 4.583$ and the signal to interference ratio is
18.66~dB at a path loss exponent of four.}

\experiment{02}{Free Space, Log-Distance and Two-Ray Ground Path Loss Models}
{To implement the free space, log-distance and two-ray ground reflection path loss models, determine
the two-ray crossover distance, and estimate the maximum cell radius for a specified receiver
sensitivity.}
{The free space path loss at 1~km and 900~MHz is 91.53~dB.}

\experiment{03}{Okumura--Hata Model and Coverage Prediction}
{To implement the Okumura--Hata model for urban, suburban and open environments, determine the
coverage radius in each, and generate a coverage map for a hexagonal base station layout designed
from the computed radius.}
{For a medium city at 900~MHz with a 30~m base station and a 1.5~m mobile, the urban median path loss
at 1~km is 126.40~dB and the mobile antenna correction is 0.0159~dB.}

\experiment{04}{Log-Normal Shadowing and Rayleigh Fading}
{To model log-normal shadowing and determine the fade margin required for a specified edge coverage
probability, to generate a Rayleigh fading envelope by the sum of sinusoids method, and to measure the
level crossing rate and average fade duration against theory.}
{The maximum Doppler shift at 60~km/h and 900~MHz is 50~Hz, and the fade margin for 90 per cent edge
coverage at $\sigma = 8$~dB is 10.25~dB.}

\experiment{05}{Diversity Techniques --- Selection, Equal Gain and Maximal Ratio Combining}
{To implement selection, equal gain and maximal ratio combining over independent Rayleigh fading
branches, and to determine the outage probability, the diversity order and the array gain as functions
of the number of branches.}
{At four branches the maximal ratio combining array gain is 6.02~dB, which is exactly
$10\log_{10}4$, and the equal gain value is 5.26~dB.}

\experiment{06}{Adaptive Equalisation of a Multipath Channel using the LMS Algorithm}
{To implement a transversal equaliser adapted by the least mean square algorithm, to compare the
converged tap weights with the Wiener solution, and to determine the reduction in intersymbol
interference and in bit error rate.}
{For the stated channel the largest eigenvalue of the autocorrelation matrix is 2.2012, giving a step
size stability limit of 0.9086, and the Wiener minimum mean square error is 0.0277.}

\experiment{07}{GSM TDMA Frame Structure and Logical Channel Mapping}
{To construct the GSM normal burst, compute the complete TDMA timing hierarchy, and map the logical
channels onto the traffic and control multiframes.}
{The TDMA frame duration is 4.6154~ms and the traffic multiframe is exactly 120~ms.}

\experiment{08}{GMSK Modulation and Spectral Analysis}
{To implement Gaussian minimum shift keying, examine the effect of the bandwidth time product on the
transmitted spectrum, and verify the constant envelope property.}
{At $BT = 0.3$ the bandwidth containing 99 per cent of the power is approximately 250~kHz, and the
envelope variation is of the order of $10^{-16}$.}

\experiment{09}{OFDM Transceiver --- Orthogonality, Cyclic Prefix and Peak to Average Power Ratio}
{To implement an OFDM transmitter and receiver, verify subcarrier orthogonality, determine the cyclic
prefix length required by a given channel, and measure the peak to average power ratio.}
{For a channel of five taps the delay spread is four samples, and the error vector magnitude stops
improving once the cyclic prefix reaches four samples.}

\experiment{10}{MIMO --- SISO, Alamouti Transmit Diversity and Spatial Multiplexing Capacity}
{To compare a single antenna link with Alamouti transmit diversity and with receive combining, and to
compute the ergodic capacity of one by one, one by two, two by one and two by two Rayleigh channels.}
{The measured diversity order is approximately 1 for the single antenna link and approximately 2 for
both Alamouti and two branch receive combining. The two by two capacity slope is close to twice the
single antenna slope.}

\experiment{11}{LTE Resource Grid, Pilot Insertion and Channel Estimation}
{To construct the LTE downlink resource grid, insert reference signals, estimate the channel by least
squares with interpolation, and measure the estimation error and the resulting error vector magnitude.}
{With a pilot spacing of six subcarriers at 15~kHz the pilots are 90~kHz apart, which is well inside the
400~kHz coherence bandwidth of a channel with 0.5~$\mu$s delay spread.}

\experiment{12}{5G NR Numerology and Waveform Comparison}
{To compute the 5G New Radio numerologies, compare their symbol and slot durations, cyclic prefix
lengths and Doppler tolerance, and relate each to its intended deployment.}
{At $\mu = 0$ the useful symbol duration is 66.67~$\mu$s and the slot duration is approximately 1~ms;
at $\mu = 1$ the slot duration is approximately 0.5~ms.}

\experiment{13}{Wireless LAN Basic Service Set --- Contention and Scheduled Access}
{To simulate a wireless LAN basic service set using the distributed coordination function, determine
the packet transmission performance as the number of stations increases, and compare it with the
scheduled access used by WiMAX.}
{At ten stations the normalised saturation throughput of the distributed coordination function is
approximately 0.76, and the analytical and simulated values agree to within one per cent.}

\experiment{14}{Mobile Ad Hoc Network with an Omnidirectional Antenna Model}
{To simulate a mobile ad hoc network of nodes with omnidirectional antennas, determine the transmission
range from a path loss model, and analyse connectivity, route length and the connectivity-based delivery
ratio against node density and transmission range.}
{At 20~dBm with a path loss exponent of 3 and a sensitivity of $-85$~dBm the transmission range is
approximately 146~m.}

\experiment{15}{Optical Fibre Link --- Attenuation, Power Budget, Digital and Analog Performance}
{To compute the attenuation spectrum of an optical fibre, establish the power budget and maximum link
length, and determine the performance of a digital link and of an analog link.}
{A bit error rate of $10^{-9}$ corresponds to a $Q$ factor of 6. The scattering component decreases
approximately with the inverse fourth power of wavelength, while the total attenuation also includes
absorption losses.}

\end{document}
